\documentclass[12pt]{article}
\usepackage{times}
\usepackage{setspace}

\setlength{\parskip}{1.5ex plus 0.5ex minus 0.5ex} % Espaciado entre párrafos
\usepackage{geometry}
\geometry{letterpaper, margin=1in}
\usepackage[
    xetex,
    pdftitle={Cache Coherence Protocol Implementation in SystemVerilog}, % Título del PDF
    pdfauthor={Montero, A.},   % Autor del PDF
    pdfsubject={Multicore Cache Coherence}, % Tema del PDF
    pdfkeywords={cache coherence, MSI, Firefly, SystemVerilog, multicore}, % Palabras clave del PDF
    pdfproducer={LaTeX with hyperref package},
    pdfcreator={pdfLatex}
]{hyperref}
\usepackage{cite}
\usepackage[backend=bibtex,style=ieee,biblabel=dot,language=spanish]{biblatex}
\addbibresource{referencias.bib}
\usepackage{times}  % Usa la fuente Times New Roman
\usepackage{svg}
\usepackage{graphicx} % Para insertar imágenes
\usepackage{float} % Para posicionar imágenes y tablas
\usepackage{caption} % Para modificar captions de tablas e imágenes
\begin{document}

%portada
% Centrar el contenido en la página
\thispagestyle{empty}
\vspace*{\fill}
\begin{center}
    \textbf{\LARGE Cache Coherence Protocol Implementation} \\[0.5cm]
    \textbf{\LARGE MSI and Firefly Protocols in SystemVerilog} \\[1.5cm]
    
    \Large Alexander Montero Vargas \\[0.3cm]
    \Large 2023166058 \\[0.8cm]
    
    \large Escuela de Ingeniería en Computadores, Tecnológico de Costa Rica \\[0.3cm]
    \large CE4302 Arquitectura de Computadores II \\[1.2cm]
    
    \large Profesor \\[1.2cm]
    
    \large 10 de mayo de 2026
\end{center}
\vspace*{\fill}

\newpage
\begin{spacing}{1.0} % Interlineado 1.0
    % DEVELOPED SYSTEM
\section{Developed System}

This section presents a detailed description of a behavioral multicore cache coherence system modeled in SystemVerilog. The system implements two cache coherence protocols (MSI and Firefly) for maintaining data consistency across multiple processing elements sharing a common memory hierarchy.

\subsection{System Architecture}

The system comprises four main components interconnected through mailbox-based message passing:

\begin{itemize}
    \item \textbf{Processing Elements (Cores):} Four independent processors that generate memory access requests (reads and writes) following traces loaded from CSV files.
    
    \item \textbf{Cache Hierarchy:} Four private L1 caches (one per core) implementing either MSI or Firefly coherence protocols. Each cache can store up to 64 lines with 32-byte line size.
    
    \item \textbf{Interconnect (Bus):} A shared bus arbiter that serializes cache requests using round-robin scheduling, broadcasts coherence events to all caches, and models bandwidth-limited memory transfers.
    
    \item \textbf{Memory:} A behavioral memory model that processes cache-induced bus requests with latency proportional to transfer size and bus bandwidth.
\end{itemize}

The system uses a class-based object-oriented design in SystemVerilog, enabling polymorphic protocol selection and modular communication through mailboxes. All timing is tracked with nanosecond precision using \texttt{\$realtime} for accurate latency measurements.

\subsection{Processing Elements and Trace Generation}

\subsubsection{Core Model}

Each \texttt{Core} class instance represents a processor that:

\begin{itemize}
    \item Generates a sequence of memory requests (PrRd: processor read, PrWr: processor write).
    \item Sends requests to its private L1 cache through a \texttt{CoreReq\_mbx} mailbox.
    \item Receives acknowledgments from the cache through a \texttt{CoreAck\_mbx} mailbox, preventing request flooding.
    \item Inserts a fixed 10 nanosecond delay between consecutive requests to simulate realistic instruction pipelining.
\end{itemize}

The core does not execute actual instructions but instead replays a pre-computed trace of memory operations, allowing reproducible and controlled behavioral simulation.

\subsubsection{Trace Generation and Workload Models}

Synthetic workloads are generated by the \texttt{workload\_gen.py} Python script, producing three distinct traffic patterns:

\begin{enumerate}
    \item \textbf{Contention Workload:} All cores repeatedly access the same cacheline, creating high coherence traffic. Includes mandatory sections where one or more cores maintain exclusive modified state (state M) to exercise write-update and write-invalidate semantics. Sequence includes minimum patterns: W, W, R, W, R to guarantee modified hits.
    
    \item \textbf{Producer-Consumer Workload:} Some cores perform repeated writes to a shared location while others read, simulating classic synchronization scenarios. This tests update propagation and invalidation-based coherence.
    
    \item \textbf{Migration Workload:} Exclusive ownership of different cachelines migrates between cores. This pattern validates state transitions (Invalid $\to$ Shared $\to$ Modified) and exercises writeback mechanisms.
\end{enumerate}

Each workload is parameterized by number of instructions (\texttt{N\_inst}) and number of cores (\texttt{Cores}). The \texttt{TraceLoader} class reads CSV files and injects requests into core queues during initialization.

\subsection{Cache Architecture}

\subsubsection{Cache Organization}

Each cache implements a direct-mapped architecture:

\begin{itemize}
    \item \textbf{Capacity:} 64 cachelines $\times$ 32 bytes = 2 KB per cache
    \item \textbf{Addressing:} 
    \begin{itemize}
        \item Bits [4:0]: Block offset (32-byte line)
        \item Bits [10:5]: Set index (64 sets)
        \item Bits [31:11]: Cache tag (21 bits)
    \end{itemize}
    \item \textbf{Line Structure:} Each line stores validity bit, tag, state, and implicit data (not modeled explicitly).
\end{itemize}

\subsubsection{Cache Coherence States}

Each cacheline exists in one of three coherence states:

\begin{itemize}
    \item \textbf{Invalid (I):} Line does not contain valid data. No copy of the line exists in this cache.
    \item \textbf{Shared (S):} Multiple caches may hold the line; the memory contains the authoritative copy. Processor reads are allowed; writes require a bus transaction.
    \item \textbf{Modified (M):} This cache exclusively owns the line; all other caches have invalid copies. This cache may write without broadcasting to the bus. Memory copy is stale.
\end{itemize}

\subsubsection{Polymorphic Protocol Selection}

The \texttt{Cache} class maintains a \texttt{ProtocolBase} handle that is instantiated with either a \texttt{ProtocolMSI} or \texttt{ProtocolFirefly} object at construction time. This design pattern (Strategy pattern) enables:

\begin{itemize}
    \item Swapping coherence protocols without modifying the Cache class.
    \item Shared metrics collection and bus interface across protocols.
    \item Extensibility for future protocol variants.
\end{itemize}

The cache issues requests through mailboxes:

\begin{itemize}
    \item \textbf{from\_core:} Receives processor requests.
    \item \textbf{to\_bus:} Sends coherence bus requests.
    \item \textbf{from\_bus:} Receives broadcast coherence events.
    \item \textbf{from\_mem:} Receives memory responses with cachelines.
\end{itemize}

\subsection{Interconnect Design}

\subsubsection{Bus Model and Arbitration}

The \texttt{Bus} class models a shared, blocking broadcast bus with the following characteristics:

\begin{itemize}
    \item \textbf{Request Collection:} Maintains per-core request queues to enable fair round-robin arbitration without modifying the shared mailbox interface.
    \item \textbf{Round-Robin Scheduling:} A pointer cycles through core queues, selecting the next non-empty queue. This ensures fairness and prevents core starvation.
    \item \textbf{Blocking Transactions:} Only one bus transaction executes at a time (no concurrent requests).
    \item \textbf{Bandwidth Modeling:} Memory transfer latency is computed as $\frac{\text{bytes}}{\text{bandwidth (bytes/ns)}}$. Default bandwidth is 1 byte/ns for shared BusRd/BusRdX transfers (32 bytes), and 0.25 byte/ns for BusUpd transfers (4 bytes).
\end{itemize}

\subsubsection{Transaction Flow and Latency}

For each arbitrated bus transaction:

\begin{enumerate}
    \item A request is dequeued and granted. The requestor's cache receives a memory response via \texttt{MemResp\_mbx}.
    \item The bus broadcasts a \texttt{BusEvent} to all caches (including requestor).
    \item All caches process the broadcast event (snooping logic).
    \item Memory processes the transaction if it is a read-like operation (BusRd or BusRdX).
    \item The bus stalls the requestor until memory responds, modeling realistic latency.
\end{enumerate}

Bus metrics tracked per transaction include:

\begin{itemize}
    \item Queue wait time: $t_{\text{grant}} - t_{\text{enqueue}}$
    \item Service time: $t_{\text{done}} - t_{\text{grant}}$
    \item Total latency: $t_{\text{done}} - t_{\text{enqueue}}$
\end{itemize}

\subsubsection{Broadcast and Snooping Mechanism}

Upon arbitration, the bus broadcasts a \texttt{BusEvent} containing:

\begin{itemize}
    \item Request type (BusRd, BusRdX, BusUpd)
    \item Address of the cacheline
    \item Source core ID
    \item Timestamp
\end{itemize}

Each cache implements a snooping task (\texttt{handle\_bus\_snoop}) that:

\begin{enumerate}
    \item Checks if it owns a copy of the requested line (valid bit, tag match).
    \item If match and request is BusRd: transitions to Shared (if not already).
    \item If match and request is BusRdX: invalidates the line (MSI) or ignores (simplified Firefly).
    \item If match and request is BusUpd: updates the line and remains Shared.
    \item Records state transitions and coherence events (invalidations, updates).
\end{enumerate}

The snooping mechanism is decoupled from the request processing logic, allowing concurrent monitoring of coherence events via the \texttt{EventMonitor} class.

\subsection{Implemented Coherence Protocols}

\subsubsection{MSI Protocol (Write-Invalidate)}

The MSI protocol uses three coherence states (Invalid, Shared, Modified) and implements write-invalidate semantics:

\paragraph{Core Request Handling:}

\begin{itemize}
    \item \textbf{PrRd (Processor Read):}
    \begin{itemize}
        \item Hit: Increment read hit counter. Line remains in current state.
        \item Miss: Issue BusRd. Memory responds with line. Cache transitions to Shared.
    \end{itemize}
    
    \item \textbf{PrWr (Processor Write):}
    \begin{itemize}
        \item Hit in M: Increment write hit counter. Write is local; no bus transaction.
        \item Hit in S: Issue BusRdX to invalidate all other copies. Cache transitions to Modified.
        \item Miss: Issue BusRdX. Memory responds. Cache transitions to Modified.
    \end{itemize}
\end{itemize}

\paragraph{Snooping Logic:}

\begin{itemize}
    \item BusRd observed: If local line is valid, mark as Shared (transition from M to S).
    \item BusRdX observed: If local line is valid, invalidate it (transition to Invalid).
    \item BusUpd observed: No action (compatibility only, not used by MSI).
\end{itemize}

\subsubsection{Firefly Protocol (Write-Update)}

The Firefly protocol implements write-update semantics, where updates to shared lines are propagated rather than invalidated:

\paragraph{Core Request Handling:}

\begin{itemize}
    \item \textbf{PrRd (Processor Read):}
    \begin{itemize}
        \item Hit: Increment read hit counter. Line remains in current state.
        \item Miss: Issue BusRd. Memory responds. Cache transitions to Shared.
    \end{itemize}
    
    \item \textbf{PrWr (Processor Write):}
    \begin{itemize}
        \item Hit in M: Increment write hit counter. Write is local; no bus transaction.
        \item Hit in S: Issue BusUpd (4-byte update transfer). Line remains Shared. All copies receive update via broadcast.
        \item Miss: Issue BusRd. Memory responds. Cache transitions to Shared, then immediately issues BusUpd if another core overwrote meanwhile (not modeled in simplified version).
    \end{itemize}
\end{itemize}

\paragraph{Snooping Logic:}

\begin{itemize}
    \item BusRd observed: Line remains Shared.
    \item BusUpd observed: Line receives update and remains Shared.
    \item BusRdX observed: Ignored (simplified Firefly does not support transition to Modified via BusRdX).
\end{itemize}

\subsubsection{Protocol Comparison}

\begin{itemize}
    \item \textbf{Write Hits on Shared Data:} MSI issues BusRdX (invalidation), Firefly issues BusUpd (update). MSI generates more bandwidth for exclusive transitions; Firefly benefits when multiple readers follow a write.
    
    \item \textbf{Bus Traffic:} MSI produces larger transactions (full cacheline invalidations). Firefly produces smaller transactions (4-byte updates).
    
    \item \textbf{Modified State Utilization:} Both use Modified state to allow write-local transactions. MSI must invalidate to reach M from S; Firefly must wait until next S-to-M transition (via subsequent miss and exclusive BusRd).
\end{itemize}

\subsection{Performance Metrics}

The system collects comprehensive metrics at multiple levels:

\subsubsection{Cache-Level Metrics}

Per-cache counters track:

\begin{itemize}
    \item \textbf{Hit Rates:} Read hits, write hits, total hits. Computed as $\text{Hit Rate} = \frac{\text{hits}}{\text{hits + misses}}$.
    \item \textbf{Miss Rates:} Read misses, write misses, total misses.
    \item \textbf{Snooping Activity:} Count of BusRd, BusRdX, BusUpd events observed.
    \item \textbf{Coherence Events:} Invalidations received, updates received, writebacks (track state changes via FSM monitor).
    \item \textbf{Bus Stalls:} Per-request stall time due to bus arbitration and memory latency. Aggregated as total, average, and maximum stall time.
\end{itemize}

\subsubsection{Bus-Level Metrics}

The bus tracks:

\begin{itemize}
    \item \textbf{Transaction Counts:} Total BusRd, BusRdX, BusUpd transactions granted.
    \item \textbf{Queue Metrics:} Total queue wait time, service time, and total latency (sum and per-transaction average).
    \item \textbf{Data Volume:} Total bytes transferred (modeled from transaction types and sizes).
    \item \textbf{Bandwidth Utilization:} Computed as $\frac{\text{bytes transferred}}{\text{simulation time}}$.
    \item \textbf{Backpressure Events:} Counts of cases where a core's request queue overflowed.
\end{itemize}

\subsubsection{Memory-Level Metrics}

Memory collects:

\begin{itemize}
    \item \textbf{Access Counts:} Total reads (BusRd), exclusive reads (BusRdX), updates (BusUpd).
    \item \textbf{Queue Statistics:} Max queue length, average queue wait, average service time, average total latency.
    \item \textbf{Transfer Volume:} Total bytes transferred and effective bandwidth.
\end{itemize}

\subsubsection{FSM Transition Monitoring}

The \texttt{EventMonitor} exports per-cache coherence state transitions to CSV format for post-simulation analysis:

\begin{itemize}
    \item Timestamp (ns)
    \item Cache ID
    \item Address and index
    \item Old state $\to$ new state
    \item Cause (e.g., ``PrRd MISS $\to$ BusRd'', ``SNOOP BusRdX'')
\end{itemize}

\subsection{Experimental Workloads}

The testbench \texttt{workload\_csv\_tb.sv} automates scenario execution and result export. Each workload generates CSV files:

\begin{itemize}
    \item Cache metrics per core
    \item Bus and memory statistics
    \item FSM transitions
    \item Summary reports
\end{itemize}

Workloads are parameterized (N\_inst, Cores) and stored with timestamps and batch identifiers for result tracking.

\end{spacing}
\printbibliography
\end{document}